\chapter{Contabilidad, colateral, gobierno e implementación}\label{cap:niif}

\begin{objetivos}
\begin{itemize}
\item Explicar la asimetría contable que crea un derivado y cómo la corrige la contabilidad de coberturas.
\item Aplicar la NIIF~9 a las dos estrategias: qué se designa, cómo y qué es inefectividad.
\item Comparar las estrategias con los criterios del Cuadro~11 del informe.
\item Cuantificar el colateral que podría exigir el banco y entender por qué es la restricción que decide.
\item Describir el plan~B y las condiciones de implementación: política, ISDA/CSA, autorizaciones y separación de funciones.
\item Releer las conclusiones del informe con todo lo aprendido.
\end{itemize}
\end{objetivos}

\begin{prerreq}Capítulos~\ref{cap:mercados}, \ref{cap:e1} y~\ref{cap:e2}.\end{prerreq}

\section{El problema contable}

Un derivado se registra a \textbf{valor razonable} (su valor de mercado) y, por regla general, sus cambios de valor van al resultado del período. La partida que se cubre, en cambio, puede medirse de otra forma: el inventario, al costo; un préstamo, al costo amortizado; una transacción futura ni siquiera está en el balance. Sin un tratamiento especial, una cobertura económicamente perfecta podría \emph{aumentar} la volatilidad de las utilidades: el derivado mueve el resultado de este trimestre y la partida cubierta lo hace en otro, o nunca.

\begin{ejemplo}{La asimetría (didáctico)}
Una empresa cubre con un NDF una compra en soles que hará dentro de seis meses. Al cierre del trimestre el sol se ha depreciado y el NDF tiene un valor negativo de 2 millones, que va a resultados. La compra futura será más barata en dólares, pero como aún no ocurrió no se registra nada. El estado de resultados muestra una pérdida de 2 millones que, económicamente, no existe.
\end{ejemplo}

\section{La contabilidad de coberturas según la NIIF 9}

La NIIF~9 (\emph{Instrumentos Financieros}, capítulo 6) permite, si se cumplen ciertos requisitos, alinear el reconocimiento del derivado y de la partida cubierta. Distingue tres tipos de relación:

\begin{definicion}[Tipos de cobertura contable]
\begin{itemize}
\item \textbf{Cobertura de valor razonable}: cubre cambios en el valor de un activo o pasivo reconocido (o de un compromiso en firme). El cambio de valor de la partida cubierta atribuible al riesgo cubierto se reconoce en resultados, junto con el del derivado.
\item \textbf{Cobertura de flujos de efectivo}: cubre la variabilidad de flujos futuros. La parte \emph{efectiva} del cambio de valor del derivado se lleva a \textbf{otro resultado integral} (ORI, una cuenta del patrimonio) y se reclasifica a resultados cuando los flujos cubiertos afectan el resultado.
\item \textbf{Cobertura de la inversión neta} en un negocio en el extranjero (no aplica aquí).
\end{itemize}
\end{definicion}

\paragraph{Requisitos} Para aplicar la contabilidad de coberturas la empresa debe: (i) \textbf{designar y documentar} la relación desde el inicio (partida cubierta, instrumento, riesgo, ratio de cobertura y método de evaluación); (ii) demostrar una \textbf{relación económica} entre partida e instrumento (sus valores se mueven en sentidos opuestos por el riesgo cubierto); (iii) que el riesgo de crédito no domine esos cambios de valor (NIIF~9, 6.4.1(c)(ii)); y (iv) usar un ratio de cobertura coherente con el que realmente gestiona. La parte \textbf{inefectiva} (la diferencia que no se compensa) va a resultados. Una transacción futura solo puede cubrirse si es \textbf{altamente probable} (6.3.3).

\begin{definicion}[Derivado hipotético]
Para medir la inefectividad de una cobertura de flujos se compara el derivado real con un \textbf{derivado hipotético}: uno imaginario que replicaría exactamente la partida cubierta (mismas fechas, montos y moneda) y que valdría cero al inicio. Lo que el derivado real cambie de más o de menos respecto del hipotético es inefectividad.
\end{definicion}

\begin{adicional}[NIIF 9 frente a NIC 39]
La norma anterior (NIC~39) exigía que la efectividad estuviera entre 80\,\% y 125\,\%. La NIIF~9 eliminó ese rango y lo reemplazó por los requisitos cualitativos anteriores, más cercanos a la gestión de riesgos real de la empresa.
\end{adicional}

\section{Estrategia 1: el CCS como cobertura de flujos; el NDF sin designar}

\begin{informe}
El CCS se designa como \textbf{cobertura de flujos de efectivo} del préstamo del BN: la parte efectiva va a ORI y pasa a resultados cuando lo hace la diferencia de cambio del préstamo. Sus fuentes de inefectividad son el cronograma, el CVA y el \emph{basis} entre monedas, que puede excluirse de la relación (NIIF~9, 6.5.16); se miden con un derivado hipotético que replica el préstamo, y se documenta que el riesgo de crédito no domina la relación. El valor razonable incluye el riesgo de crédito propio (DVA; NIIF~13, ¶42).

El NDF, en cambio, \textbf{no se designa}: la diferencia de cambio de las cuentas por pagar va a resultados (NIC~21) y el NDF la compensa en el mismo estado. No hace falta contabilidad de coberturas porque ambas partidas ya se miden en resultados.
\end{informe}

\begin{ejemplo}{Una caída instantánea de 10\,\% del TC}
Con el TC medio pasando de 3.4379 a 3.0941:
\begin{itemize}
\item El saldo del préstamo del BN (S/ 3,765.6 millones) aumenta en dólares: $3{,}765.6/3.0941-3{,}765.6/3.4379=1{,}217.0-1{,}095.3=121.7$ millones. Es una pérdida por diferencia de cambio en resultados.
\item El CCS pasa a ser un activo de US\$ 123.7 millones. Esa ganancia se registra en ORI.
\item Se reclasifican de ORI a resultados 121.7 millones, que compensan la pérdida cambiaria del préstamo.
\item El resultado del período queda prácticamente sin efecto; en ORI quedan unos 2.0 millones, que corresponden a la parte del valor del CCS ligada a los intereses futuros (el cupón del BN supera la curva) y se reclasificarán cuando esos intereses afecten el resultado.
\end{itemize}
Al mismo tiempo, las cuentas por pagar en soles pierden y el NDF (activo de 24.0 millones) lo compensa directamente en resultados.
\end{ejemplo}

\section{Estrategia 2: cobertura de valor razonable de un inventario que rota}

\begin{informe}
Las tres capas suman el 80\,\% del inventario que \emph{existe hoy} y se designan como \textbf{cobertura de valor razonable} de ese inventario, por su \emph{riesgo de precio total} y en ratio 1:1. No se designan como cobertura de un \emph{componente} WTI, porque el crudo se compra referido al Brent. Al venderse un lote, la relación se discontinúa (NIIF~9, 6.5.6) y el derivado se redesigna sobre el lote que lo reemplaza; su valor a esa fecha y la base Brent--WTI son inefectividad. En el collar se excluye el valor temporal como costo de la cobertura (6.5.15).
\end{informe}

Tres ideas se combinan en ese párrafo:
\begin{itemize}
\item \textbf{Riesgo total, no componente.} La NIIF~9 permite cubrir solo un \enquote{componente} del precio (por ejemplo, la parte WTI) si es identificable y medible. Como el crudo de PETROPERÚ sigue al Brent, se designa el riesgo de precio completo y se acepta que la base genere inefectividad.
\item \textbf{Discontinuar y redesignar.} El inventario rota cada mes: el lote cubierto se vende. Cada vez, la relación termina y el derivado (que sigue vivo) se asigna al lote nuevo. Como el derivado ya no vale cero en ese momento, su valor es fuente de inefectividad.
\item \textbf{Valor temporal del collar.} Los cambios del valor tiempo de las opciones (por volatilidad o por el paso del tiempo) no tienen contrapartida en el inventario. Excluirlos de la relación y tratarlos como \enquote{costo de la cobertura} evita esa inefectividad.
\end{itemize}

\begin{cuidado}[escenario económico no es pérdida contable]
El informe advierte que su Cuadro~8 es \emph{económico}. La pérdida contable del inventario depende de su costo y de su valor neto de realización (NIC~2): un inventario comprado a 57.42 no registra pérdida contable si el WTI cae de 92.60 a 80, aunque económicamente haya perdido valor frente a hoy.
\end{cuidado}

\section{Comparación de las estrategias}

\begin{center}\small
\captionof{table}{Comparación de las estrategias (Cuadro~11 del informe), US\$ millones.}\label{tab:c10}
\begin{tabularx}{\linewidth}{@{}>{\raggedright\arraybackslash}p{0.2\linewidth}>{\raggedright\arraybackslash}X>{\raggedright\arraybackslash}X>{\raggedright\arraybackslash}X@{}}
\toprule
Criterio & E1: CCS + NDF & E2: swap + collar & Tasa (no hoy)\\
\midrule
Exposición cubierta & 96\,\% de la posición en S/ & 2.32 MMbl, 80\,\% del inventario (214.9) & Tasa base del CESCE (722.2)\\
Reducción del riesgo & TC $-10\,\%$: de (152.0) a (6.1) & WTI $-30\,\%$: de (80.6) a (34.0) & Baja: no cubre el spread\\
Costo inicial & Cero; CVA de 11 pb en la tasa & Cero & Cero (inicio diferido)\\
Costo económico & 0.24\,\% anual (NDF) & Swap: 3.7 frente al spot (nulo frente al futuro) & Pagos si la tasa baja\\
Colateral con umbral cero & 121.6 si el TC sube 10\,\% & 51.8 si el WTI sube 30\,\% & Bajo\\
NIIF 9 & CCS: flujos (ORI); NDF: resultados & Valor razonable; se redesigna al rotar & Flujos, solo si es altamente probable\\
\bottomrule
\end{tabularx}
\end{center}

La fila nueva respecto de los capítulos anteriores es la del colateral. El informe la resume en su frase central: \enquote{lo que puede impedir ejecutarlas no es su costo, sino el colateral que pediría el banco}.

\section{Liquidez: la restricción que decide}

Los derivados OTC no tienen márgenes de bolsa, pero un banco pedirá colateral a una contraparte B-/Caa1. El informe supone el caso más exigente, \textbf{umbral cero}, y mide cuánto podría tener que entregar PETROPERÚ.

\begin{center}\small
\captionof{table}{Colateral exigible con umbral cero, US\$ millones (Cuadro~13 del informe). Déficit: colateral menos líneas libres (22.2).}\label{tab:colateral-libro}
\begin{tabular}{@{}lrrrr@{}}
\toprule
& E1 & E2 & Colat. & Déficit\\
\midrule
TC $+10\,\%$ & 121.6 & 0.0 & 121.6 & 99.4\\
WTI $+30\,\%$ & 0.4 & 51.8 & 52.2 & 30.0\\
Ambos a la vez & 121.6 & 51.8 & 173.4 & 151.2\\
PFE 95\,\% & 108.3 & 97.5 & 168.9 & 146.7\\
PFE 99\,\% & 151.4 & 157.8 & 253.6 & 231.4\\
\bottomrule
\end{tabular}
\end{center}

Cómo se lee:
\begin{itemize}
\item \textbf{Choques instantáneos.} Si el TC sube 10\,\%, el CCS y el NDF son pasivos por $101.9+19.7=121.6$ millones. Si el WTI sube 30\,\%, el swap y el collar lo son por 51.8. A la vez, 173.4.
\item \textbf{PFE.} Pero el CCS vive 27 meses: lo relevante no es un choque de hoy sino la exposición potencial futura en toda su vida (cap.~\ref{cap:mercados}). Para la E1 se toma el pico de la PFE del CCS más el NDF en el cuantil a 91 días (151.4 al 99\,\%); para la E2, el programa rodante en el cuantil a su plazo (157.8).
\item \textbf{Agregación.} E1 y E2 se suman con la correlación TC--WTI de 0.35 (cap.~\ref{cap:riesgo}): $\sqrt{151.4^2+157.8^2+2\times0.35\times151.4\times157.8}\approx253.6$.
\end{itemize}

Frente a esos 253.6 millones, PETROPERÚ tiene US\$ 22.2 millones de líneas libres ($376.9-354.6$). El préstamo puente de US\$ 475 millones que obtuvo con apoyo del Gobierno tiene otro destino (normalizar el capital de trabajo).

\begin{ficha}{umbral necesario del CSA}
\paso{1. Ecuación}
\[ \text{Umbral}=\text{PFE}_{99\,\%}-(1-\text{colchón})\times\text{Líneas libres}. \]
\paso{2. Símbolos} PFE$_{99\,\%}$: colateral conjunto en el escenario del 99\,\%. Colchón: fracción de las líneas libres que se reserva para la operación (25\,\%).
\paso{3. Intuición} El banco debe tolerar sin pedir colateral todo lo que la empresa no puede pagar con sus líneas, dejando una reserva para operar.
\paso{4. Construcción} Si el colateral exigible en el escenario malo es la PFE y la empresa puede usar a lo sumo el 75\,\% de sus líneas libres, el resto debe cubrirlo el umbral.
\paso{5. Supuestos} PFE lognormal con la volatilidad realizada; tres bancos con exposiciones iguales.
\paso{6. Ejemplo (del informe)} $253.6-0.75\times22.2=253.6-16.7=236.9\to237$ millones; con tres bancos, unos 79 millones cada uno.
\paso{7. Interpretación} La ejecución exige un umbral total de al menos US\$ 237 millones o la garantía del Estado. \enquote{Sin ello, la cobertura agravaría el riesgo de liquidez que busca mitigar.}
\paso{8. Limitaciones} El 99\,\% no es el peor caso; la volatilidad puede subir; un banco puede reducir el umbral si baja la calificación.
\end{ficha}

\subsection{Plan B si el banco no concede el umbral}

\begin{enumerate}
\item[(a)] \textbf{Escalonar}: contratar primero el 50\,\% del CCS. El colateral ante un TC $+10\,\%$ baja de 121.6 a 70.7 millones.
\item[(b)] \textbf{Comprar opciones}: una opción comprada solo exige prima y nunca colateral. La put de WTI sobre el 80\,\% cuesta US\$ 7.2 millones, por lo que se limitaría al tramo más expuesto.
\item[(c)] \textbf{Mejorar el CSA}: umbral escalonado según la calificación y reajuste semanal del colateral en lugar de diario.
\item[(d)] \textbf{Garantía del Estado} sobre la línea de derivados.
\end{enumerate}

\section{Riesgo de tasa: no se recomienda cubrirlo hoy}

La deuda de largo plazo ya es fija; pasarla a variable añadiría riesgo sin bajar el costo de crédito, que lo domina el spread. Queda como opción el swap de inicio diferido para la refinanciación del CESCE (4.91\,\% con pago anual; cap.~\ref{cap:swaps}), solo cuando esa refinanciación sea altamente probable (NIIF~9, 6.3.3). Hoy no lo es, por el incumplimiento del covenant.

\section{Condiciones de implementación}

Un derivado mal gobernado puede convertirse en especulación. El informe propone estas condiciones:

\begin{enumerate}
\item \textbf{Política de Cobertura aprobada por el Directorio}, con límites (70\,\%--100\,\% de la posición en soles y 50\,\%--80\,\% del inventario de crudo), prohibición expresa de especular y métricas mensuales. Los límites evitan tanto la falta de cobertura como la sobrecobertura. Incluye límites por contraparte y de pérdida.
\item \textbf{Contratos ISDA con anexo CSA} y un umbral de al menos US\$ 237 millones, o el plan~B, con período de reposición del colateral y monto mínimo de transferencia. Como el covenant del CESCE ya está incumplido, el incumplimiento cruzado (§5(a)(vi)) se limita a \emph{cross-acceleration}, con umbral alto y dispensa expresa para el CESCE; sin ella, no se firma (cap.~\ref{cap:mercados}). Se negocia además que la rebaja de calificación, la opinión con salvedades o la incertidumbre de empresa en marcha no causen el cierre inmediato sin un período de cura, y una cláusula de prepago: si el BN prepaga, el CCS se cancela a valor de mercado.
\item \textbf{Opinión legal y autorizaciones}: antes de pactar, la Gerencia Legal opina y se obtiene la autorización que exija el régimen de endeudamiento y tesorería de una empresa del Estado. Se evalúa si los compromisos contingentes que el D.U. N.° 003-2026 autoriza al MINEM (hasta US\$ 2,000 millones) pueden respaldar las líneas de derivados.
\item \textbf{Documentación y revelación}: las coberturas se documentan según la NIIF~9 y se revelan según la NIIF~7 (¶21A--24G); los swaps con Citibank entran a la política. Con insumos aproximados (\emph{proxy}) y el CVA, el nivel de jerarquía de valor razonable (2 o 3; NIIF~13, ¶73) se determina por instrumento.
\item \textbf{Separación de funciones}: Finanzas contrata (\emph{front office}), Riesgos controla límites, aprueba los modelos y verifica de forma independiente los precios de los bancos (\emph{middle office}), y Contabilidad registra (\emph{back office}).
\end{enumerate}

\begin{informe}[cuándo revisar y qué vigilar]
Se revisa la decisión si el BN prepaga su préstamo; si el CVA, el FVA o el \emph{basis} superan el beneficio de estabilizar la caja; o si el marcador de venta resulta ser el Brent (la E2 pasaría a swaps de Brent). Cada mes se vigilan el colateral exigido frente a las líneas libres, la posición neta en soles frente a lo cubierto, el diferencial Brent--WTI y la calificación crediticia, que activa los umbrales del CSA.
\end{informe}

\section{Por qué todo esto vale la pena}

Volvamos a la pregunta del capítulo~\ref{cap:caso}. En un mundo perfecto, cubrirse no crearía valor. PETROPERÚ está lejos de ese mundo: pierde dinero, su capital de trabajo es negativo en US\$ 2,296.4 millones, usa el 94\,\% de sus líneas, el mercado le exige 9.7\,\% en dólares y el auditor duda de su continuidad. En esas condiciones:
\begin{itemize}
\item una pérdida cambiaria de US\$ 152 millones o una caída del crudo de US\$ 80 millones no se absorbe con caja propia: obliga a buscar financiamiento caro o a dejar de operar (Froot et al., 1993);
\item el objetivo no es ganar con los precios sino eliminar los escenarios extremos que acercan a la empresa a la insolvencia (Stulz, 1996);
\item los costos de la cobertura (0.24\,\% anual en el NDF, 3.7 millones de backwardation frente al spot, 11 pb de CVA) son conocidos y acotados, frente a pérdidas potenciales mucho mayores e inciertas.
\end{itemize}
Pero el caso también muestra el límite de esa lógica: si la cobertura exige más colateral del que la empresa puede entregar, convierte un riesgo de precio en un riesgo de liquidez. Por eso la recomendación es condicional.

\begin{adicional}[limitaciones del análisis]
El propio informe declara que sus cifras son indicativas: (i) el cronograma del BN es inferido (tasa del CCS entre 5.82\,\% y 6.01\,\%); (ii) la posición en soles distinta del BN y el inventario son de diciembre de 2025: si cambiaron, cambian los nocionales, no el diseño; (iii) la curva soberana y el OVX sustituyen a la curva swap en soles y a la superficie de volatilidad, que no son públicas; (iv) el CVA es un mínimo (sin FVA, KVA ni riesgo \emph{wrong-way}); (v) los escenarios de la E2 suponen que el inventario sigue al WTI y sobrestiman la reducción del riesgo; (vi) el VaR y la PFE suponen lognormalidad y volatilidad constante. Antes de pactar se requieren el cronograma contractual, la posición actualizada y cotizaciones ejecutables de al menos tres bancos (no BlueGamma ni Yahoo Finance).
\end{adicional}

\section{Las conclusiones del informe, releídas}

\begin{enumerate}
\item \textbf{Decisión}: aprobar las dos estrategias. PETROPERÚ asume hoy, sin cobertura, una deuda neta de S/ 4,704.0 millones y un inventario expuesto a un WTI que cayó 19.9\,\% en 2025 y hoy cotiza a US\$ 92.60/bl. Con utilidad bruta casi nula y liquidez restringida, una mala racha compromete la continuidad operativa. \emph{(Capítulos 1, 3 y 10.)}
\item \textbf{Estrategia 1}: el CCS al 5.98\,\% \emph{all-in} (5.93\,\% con el spread de agosto de 2026) más NDF cubre el 96\,\% de la posición en soles; ante una caída de 10\,\% del TC, la pérdida baja de US\$ 152.0 a 6.1 millones, sin pago inicial. \emph{(Capítulos 2, 4, 5, 7 y 11.)}
\item \textbf{Estrategia 2}: un programa rodante sobre el 80\,\% del inventario, mitad swap y mitad collar de costo cero, reduce la pérdida ante una caída de 30\,\% del WTI de US\$ 80.6 a 34.0 millones, sin prima. El margen de refinación queda para una fase~2; la tasa no se cubre hasta que la refinanciación del CESCE sea altamente probable. \emph{(Capítulos 6, 8, 9 y 12.)}
\item \textbf{Condición}: se ejecuta solo con un umbral CSA de al menos US\$ 237 millones o con la garantía del Estado, porque al 99\,\% el colateral podría llegar a US\$ 253.6 millones en la vida del CCS; y con dispensa del incumplimiento cruzado por el CESCE. Sin esa condición, plan~B o se pospone. \emph{(Capítulos 4 y 13.)}
\item \textbf{Próximos pasos}: presentar la Política de Cobertura al Directorio; negociar ISDA/CSA con al menos tres bancos; obtener la opinión legal y el respaldo del Estado (MINEM, D.U. N.° 003-2026); y ejecutar primero el CCS, que cubre el riesgo de mayor impacto.
\end{enumerate}

\section{Guía para releer el informe}

Al terminar este libro, cada término del informe debería remitir a una idea precisa:

\begin{center}\small
\begin{tabularx}{\linewidth}{@{}>{\raggedright\arraybackslash}p{0.5\linewidth}X@{}}
\toprule
Término del informe & Dónde se explica\\
\midrule
Moneda funcional, empresa en marcha, calificación & Cap.~1\\
Curva cupón cero, bootstrap, SOFR, tasa efectiva, cuota francesa con gracia & Cap.~2\\
Volatilidad, asimetría $\pm10\,\%$, duración, VaR, correlación & Cap.~3\\
OTC, márgenes, ISDA, CSA, umbral, EPE, PFE, CVA, cross-acceleration & Cap.~4\\
NDF, paridad cubierta, pips, TC de compra & Cap.~5\\
Backwardation, rendimiento de conveniencia, base, $h^*$ & Cap.~6\\
CCS, swap a precio promedio, swap de inicio diferido & Cap.~7\\
Put, call, collar de costo cero, paridad & Cap.~8\\
Modelo de Black, volatilidad implícita, skew, delta, vega & Cap.~9\\
Matriz de riesgos, margen de refinación (\emph{crack}) & Cap.~10\\
PCC, PND, PCG & Cap.~11\\
Saldo estructural, programa rodante, capas & Cap.~12\\
NIIF 9, ORI, derivado hipotético, colateral, plan~B & Cap.~13\\
\bottomrule
\end{tabularx}
\end{center}

\begin{resumen}
\begin{itemize}
\item Sin contabilidad de coberturas, los derivados pueden aumentar la volatilidad contable aunque reduzcan la económica.
\item E1: el CCS es cobertura de flujos de efectivo (ORI, derivado hipotético, \emph{basis} excluible); el NDF no se designa porque ambas partidas ya van a resultados.
\item E2: cobertura de valor razonable del inventario existente, por su riesgo total; se discontinúa y redesigna al rotar; el valor temporal del collar se excluye.
\item El colateral es la restricción que decide: PFE conjunta al 99\,\% de 253.6 millones frente a 22.2 de líneas libres. La ejecución exige un umbral de al menos 237 millones (79 por banco) o la garantía del Estado; si no, plan~B.
\item La implementación exige política con límites, ISDA/CSA con \emph{cross-acceleration} y período de cura, autorizaciones legales, documentación NIIF~9 y separación de funciones.
\end{itemize}
\end{resumen}

\section*{Ejercicios}
\addcontentsline{toc}{section}{Ejercicios}

\ejercicio[conceptual] Explique con un ejemplo propio por qué un derivado que reduce el riesgo económico puede aumentar la volatilidad de las utilidades si no se aplica contabilidad de coberturas.

\ejercicio[cálculo] Con una subida instantánea de 10\,\% del TC medio (de 3.4379 a 3.7817), calcule la ganancia cambiaria del préstamo del BN en dólares y describa el registro bajo cobertura de flujos de efectivo, sabiendo que el CCS pasa a ser un pasivo de US\$ 101.9 millones.

\ejercicio[conceptual] ¿Por qué la cobertura del inventario se designa como de valor razonable, la del préstamo como de flujos de efectivo y el NDF no se designa?

\ejercicio[conceptual] ¿Qué riesgo busca evitar el límite superior de la política (100\,\% de la posición en soles, 80\,\% del inventario)?

\ejercicio[cálculo] Si las líneas libres fueran de US\$ 60 millones y la PFE conjunta al 99\,\% siguiera en 253.6, ¿qué umbral total exigiría la regla del informe? ¿Cuánto por banco con tres bancos?

\ejercicio[conceptual] ¿Por qué es importante que quien contrata los derivados no sea quien mide los límites ni quien los registra?

\ejercicio[integrador] Redacte, en no más de diez líneas, la recomendación a la Gerencia General explicando por qué PETROPERÚ debería implementar ambas estrategias, mencionando beneficio, costo, condición de ejecución y principal riesgo residual de cada una.
